\chapter{إحصاء أصفار دالة زيتا والمصفوفات العشوائية}
\label{ch:zero-statistics-random-matrices}

\section*{نطاق الفصل وحالته}

ينتقل هذا الفصل من عدّ أصفار دالة زيتا إلى دراسة بنيتها المحلية بعد إزالة تغير الكثافة مع الارتفاع. النواة الصارمة هي صيغة مونتغمري للارتباط الثنائي ضمن مجال دعم محدود وتحت فرضية ريمان في الصياغة التاريخية المعتمدة هنا. أما صيغة GUE الكاملة، وتوزيع الفواصل المتتالية، وإحصاءات المستويات الأعلى فتبقى حدسية أو جبهة مفتوحة بحسب السياق.

\noindent\textbf{حالة الفصل:}
\textenglish{OWNER-ADOPTED / ACTIVE / CITABLE}. اجتاز الفصل البناء الكامل والمراجعة المستقلة بعد التأليف بحكم \textenglish{PASS / 0 BLOCKERS}، واعتمده المالك صراحة في 26 يوليو 2026. تبقى التصنيفات العلمية لكل معرّف كما هي: المبرهنات المستشهد بها، والحدسيات، والأدلة العددية، والمبادئ المنهجية لا تختلط فيما بينها.

\subsection*{حراس النطاق}
\begin{itemize}
  \item صيغة ريمان--فون مانغولت لا تسجل هنا كمبرهنة جديدة؛ مصدرها الداخلي هو \texttt{ANT-THM-06-06} في الفصل السادس.
  \item مبرهنة مونتغمري المعروضة هنا مشروطة بفرضية ريمان ومقيدة بدعم فورييه داخل \((-1,1)\).
  \item حد GUE مبرهنة في نموذج المصفوفات العشوائية، لا مبرهنة عن أصفار زيتا.
  \item حسابات Odlyzko دليل عددي منته، وليست برهانًا على RH أو على حدسية GUE.
  \item لا يساوى الارتباط الثنائي بتوزيع أقرب جار أو بتباين العدد أو بإحصاءات المستويات الأعلى.
\end{itemize}

\subsection*{نواتج التعلم}
بنهاية الفصل يستطيع القارئ أن:
\begin{enumerate}
  \item يستخرج كثافة الأصفار المتوسطة ومتوسط التباعد من صيغة عد الأصفار.
  \item يطبق التطبيع المحلي الذي يجعل متوسط التباعد قريبًا من واحد.
  \item يميز بين دالة مونتغمري الموزونة وصيغة دوال الاختبار للأزواج غير القطرية.
  \item يحدد بدقة موضع شرط RH وقيد دعم التحويل الفوري.
  \item يفسر ظهور نواة الجيب في GUE دون تحويل النموذج إلى برهان عن زيتا.
  \item يفرق بين الدليل العددي والنتيجة التقاربية والحدسية.
  \item يقرأ مبدأ Katz--Sarnak مع حارس يمنع النقل الآلي بين العائلات.
\end{enumerate}

\section{عد الأصفار والمقياس المحلي}

\begin{definition}[دالة عد الأصفار والتطبيع المحلي]
\label{def:ch23-zero-count-local-scale}
\resultid{ANT-DEF-23-01}
\provedhere
نكتب
\[
N(T)=\#\{\rho=\beta+i\gamma:0<\gamma\le T\},
\]
مع العد بالتعدد. وبالرجوع إلى مبرهنة ريمان--فون مانغولت المثبتة في الفصل السادس، \texttt{ANT-THM-06-06}، لدينا
\[
N(T)=\frac{T}{2\pi}\log\frac{T}{2\pi}-\frac{T}{2\pi}+O(\log T).
\]
ومن مشتقة الحد الرئيس تكون الكثافة المتوسطة قرب الارتفاع \(T\)
\[
\rho_T=\frac{1}{2\pi}\log\frac{T}{2\pi},
\]
ومتوسط التباعد المحلي
\[
\Delta_T=\rho_T^{-1}=\frac{2\pi}{\log(T/(2\pi))}.
\]
ولهذا يطبع فرق ارتفاعين بالصورة
\[
u=(\gamma-\gamma')\frac{\log T}{2\pi},
\]
مع جواز استبدال \(\log T\) بـ\(\log(T/(2\pi))\) في الحدود المحلية التي لا تتأثر إلا بخطأ نسبي من رتبة \(1/\log T\).
\end{definition}

هذا التطبيع ليس تجميلًا شكليًا. فالفواصل الخام تنكمش مع الارتفاع، بينما يجعل المتغير \(u\) متوسط التباعد قريبًا من واحد، وهو المقياس الذي تقارن عليه أصفار زيتا بطيف المصفوفات العشوائية.

\section{الارتباط الثنائي واتفاقية فورييه}

\begin{definition}[إحصاء الأزواج ودالة مونتغمري]
\label{def:ch23-pair-correlation}
\resultid{ANT-DEF-23-02}
\provedhere
نعتمد اتفاقية فورييه
\[
\widehat f(\alpha)=\int_{\mathbb R}f(u)e^{-2\pi i\alpha u}\,du.
\]
ونضع
\[
w(u)=\frac{4}{4+u^2}.
\]
تحت RH نكتب الأصفار \(\rho=1/2+i\gamma\)، ونعرّف
\[
F(\alpha,T)=
\left(\frac{T}{2\pi}\log T\right)^{-1}
\sum_{0<\gamma,\gamma'\le T}
T^{i\alpha(\gamma-\gamma')}w(\gamma-\gamma').
\]
الأزواج هنا مرتبة والقطر \(\gamma=\gamma'\) داخل المجموع. أما إحصاء دالة الاختبار غير القطري فيأخذ الصورة
\[
\frac{1}{N(T)}
\sum_{\substack{0<\gamma,\gamma'\le T\\ \gamma\ne\gamma'}}
f\!\left((\gamma-\gamma')\frac{\log T}{2\pi}\right).
\]
لا تخلط الصيغتان دون حساب التحويل الفوري والقطر والوزن.
\end{definition}

\section{مبرهنة مونتغمري ومجال الدعم}

\begin{theorem}[مونتغمري: الجزء المثبت من الارتباط الثنائي]
\label{thm:ch23-montgomery}
\resultid{ANT-THM-23-01}
\citedresult[\cite{montgomery1973pair}]
نفترض فرضية ريمان. للصيغة الموزونة السابقة، وعندما \(0\le\alpha<1\)،
\[
F(\alpha,T)
=
T^{-2\alpha}(\log T+O(1))+\alpha+o(1),
\]
بصورة منتظمة على المجالات المغلقة داخل هذا المجال. وبصيغة دوال الاختبار، يطابق حد الأزواج غير القطرية كثافة نواة الجيب عندما
\[
\operatorname{supp}\widehat f\subset(-1,1).
\]
\end{theorem}

القيد على الدعم هو الحد الفاصل بين المبرهنة والحدسية. لا يسمح النص السابق بتمديد الصيغة تلقائيًا إلى جميع \(\alpha\)، ولا إلى كل دالة اختبار. وتوفر إعادة الصياغة المنشورة حديثًا في \cite{baluyot2024unconditional} سياقًا تدقيقيًا معاصرًا، لكنها لا تلغي الإسناد التاريخي إلى مونتغمري ولا تغير تصنيف الصيغة المعتمدة هنا.

الحد \(T^{-2\alpha}\log T\) يمثل أثر القطر والبنية القريبة منه في الصيغة الموزونة. وعندما \(0<\alpha<1\) يكون \(F(\alpha,T)\to\alpha\)، ومن ثم يقترب \(1-F(\alpha,T)\) من \(1-\alpha\)، وهو الجزء المثلثي \((1-|\alpha|)_+\) على هذا المجال.

\section{حدسية الارتباط الثنائي الكاملة}

\begin{conjecture}[حدسية مونتغمري--GUE للارتباط الثنائي]
\label{conj:ch23-pair-correlation}
\resultid{ANT-CONJ-23-01}
للأزواج غير القطرية وبعد جعل متوسط التباعد واحدًا، يتوقع أنه لكل دالة اختبار مناسبة
\[
\frac{1}{N(T)}
\sum_{\substack{0<\gamma,\gamma'\le T\\ \gamma\ne\gamma'}}
f\!\left((\gamma-\gamma')\frac{\log T}{2\pi}\right)
\longrightarrow
\int_{\mathbb R}f(u)
\left[1-\left(\frac{\sin\pi u}{\pi u}\right)^2\right]du.
\]
\end{conjecture}

إذا وضعنا
\[
K(u)=\frac{\sin\pi u}{\pi u},
\qquad K(0)=1,
\]
فالكثافة المتوقعة هي
\[
R_2(u)=1-K(u)^2.
\]
وعند اتفاقية فورييه المعتمدة لدينا
\[
\widehat{K^2}(\alpha)=(1-|\alpha|)_+.
\]
هذه الهوية تفسر لماذا يرى برهان مونتغمري الجزء ذي الدعم المحدود، لكنها لا تزيل الحاجز خارج \((-1,1)\).

\section{نواة الجيب في نموذج GUE}

\begin{theorem}[حد GUE المحلي ونواة الجيب]
\label{thm:ch23-gue-sine-kernel}
\resultid{ANT-THM-23-02}
\citedresult[\cite{mehta2004random}]
في حد الحجم الكبير للمجموعة الوحدوية الهرميتية، وبعد التطبيع المحلي بحيث تصبح الكثافة واحدًا، تتقارب نواة الارتباط إلى
\[
K(u-v)=\frac{\sin\pi(u-v)}{\pi(u-v)}.
\]
ومن ثم تكون كثافة الارتباط الثنائي غير القطري
\[
R_2(u)=1-\left(\frac{\sin\pi u}{\pi u}\right)^2.
\]
\end{theorem}

هذه مبرهنة عن نموذج مصفوفي. التشابه مع صيغة أصفار زيتا عميق ومثمر، لكنه ليس تطابقًا مثبتًا بين الكائنين. الفارق المنطقي يجب أن يبقى ظاهرًا: في GUE تأتي النواة من بنية محددات الارتباط، أما في زيتا فالصيغة الكاملة ما تزال حدسية.

\section{الدليل العددي: ما يثبته وما لا يثبته}

\begin{remark}[دليل Odlyzko العددي]
\label{rem:ch23-odlyzko}
\resultid{ANT-EVID-23-01}
\citedresult[\cite{odlyzko1987spacings}]
قارن Odlyzko إحصاءات الفواصل لأول \(10^5\) صفر، ثم لكتلة من \(10^5\) صفر تبدأ عند الفهرس \(10^{12}+1\)، مع دقة عددية معلنة تقارب \(10^{-8}\). أظهرت البيانات توافقًا أقوى مع تنبؤات GUE عند الارتفاع الأكبر.
\end{remark}

التصنيف الصحيح هو
\[
\textenglish{NUMERICAL-EVIDENCE / FINITE-VERIFIED}.
\]
حتى عينة ضخمة لا تختبر إلا عددًا منتهيًا من الأصفار، ولا تستبعد سلوكًا مختلفًا لاحقًا، ولا تثبت أن جميع الأصفار تقع على الخط الحرج.

\section{الفواصل وتباين العدد وإحصاءات المستويات}

\begin{definition}[إحصاءات محلية متمايزة]
\label{def:ch23-local-statistics}
\resultid{ANT-DEF-23-03}
\provedhere
لتكن \(\widetilde\gamma_n\) الأصفار بعد التطبيع المحلي.
\begin{enumerate}
  \item \textbf{فاصل أقرب جار:} توزيع \(\widetilde\gamma_{n+1}-\widetilde\gamma_n\).
  \item \textbf{الارتباط الثنائي:} متوسط على جميع الأزواج المختلفة بدالة اختبار.
  \item \textbf{تباين العدد:} إذا كان \(N_I\) عدد النقاط في نافذة طولها \(L\)، فننظر إلى \(\operatorname{Var}(N_I)\).
  \item \textbf{إحصاء المستوى \(n\):} متوسطات على \(n\)-توائم متميزة من النقاط.
\end{enumerate}
هذه الكائنات مترابطة لكنها غير مترادفة.
\end{definition}

معرفة الارتباط الثنائي الكامل قد تحدد بعض المتوسطات التربيعية وتباينات العد عبر التكامل، لكنها لا تحدد تلقائيًا قانون أقرب جار؛ فالأخير يعتمد على غياب نقاط أخرى بين الجارين، وهي معلومة تتداخل فيها جميع المستويات.

\section{عائلات دوال \(L\) وأنواع التناظر}

\begin{remark}[مبدأ Katz--Sarnak]
\label{rem:ch23-katz-sarnak}
\resultid{ANT-PRIN-23-01}
\citedresult[\cite{katzsarnak1999random}]
عند دراسة عائلة مضبوطة من دوال \(L\)، لا يكفي قول إن الأصفار «تشبه المصفوفات العشوائية». يجب تحديد العائلة، ومقياس الموصل، وحد الكثافة، ونوع التناظر المتوقع: وحدوي أو تعامدي أو سمبلكتي.
\end{remark}

تقدم أعمال Katz--Sarnak مبرهنات قوية في عائلات هندسية فوق الحقول المنتهية وتبني قاموس التناظر. أما نقل هذا القاموس إلى عائلات دوال \(L\) الكلاسيكية فهو مبدأ تنبؤي مدعوم بنتائج جزئية، لا عملية نقل آلي للثوابت أو مجالات الدعم. وتوفر نتيجة Rudnick--Sarnak \cite{rudnicksarnak1996zeros} سياقًا لإحصاءات المستويات في دوال \(L\) الرئيسية ضمن قيود دعم مناسبة.

\section{حدود الاستدلال}

\begin{remark}[حدود ما تسمح به البيانات والإحصاءات]
\label{rem:ch23-inference-limits}
\resultid{ANT-PRIN-23-02}
\begin{center}
\small\textenglish{METHODOLOGICAL-PRINCIPLE}\par
\small\textenglish{INFERENCE-GUARDED}
\end{center}
توافق إحصاء pair correlation مع GUE لا يثبت RH؛ فقد صيغت مبرهنة مونتغمري التاريخية هنا تحت RH أصلًا، ولا ينتج من موافقة متوسط زوجي تحديد موضع كل صفر. كذلك لا يستعيد الإحصاء الثنائي وحده قانون الفواصل المتتالية أو كل إحصاءات \(n\)-level.
\end{remark}

هذا العنصر ليس مبرهنة جديدة ولا يحمل تصنيف \textenglish{PROVED-HERE}. وظيفته منع قفزات استدلالية شائعة: من متوسط إلى كل نقطة، ومن إحصاء ثنائي إلى البنية الكاملة، ومن دليل عددي منته إلى عبارة لا نهائية.

\section{الجبهة المفتوحة}

\begin{openproblem}[GUE وإحصاءات المستويات]
\label{open:ch23-full-gue}
\resultid{ANT-OPEN-23-01}
\openresult
إثبات حدسية الارتباط الثنائي بلا قيد الدعم، ثم إثبات إحصاءات المستويات الأعلى وقوانين الفواصل بالقوة المتوقعة لأصفار زيتا وعائلات دوال \(L\)، يبقى من الجبهات المركزية. ويجب في كل عائلة تحديد التناظر والموصل والتطبيع قبل صياغة الحدسية.
\end{openproblem}

\section*{خلاصة الفصل}

ينشأ المقياس المحلي من صيغة عد الأصفار، ثم يكشف تحويل فورييه أن مبرهنة مونتغمري ترى جزءًا محدود الدعم من بنية نواة الجيب. في نموذج GUE تكون النواة مبرهنة، بينما تبقى الصيغة الكاملة لأصفار زيتا حدسية. بيانات Odlyzko تقدم دليلًا عدديًا قويًا لكنه منته. وأخيرًا، لا يجوز الخلط بين الارتباط الثنائي والفواصل المتتالية أو إحصاءات المستويات، ولا بين قاموس Katz--Sarnak والبرهان في كل عائلة كلاسيكية.

\section*{تمارين}
\begin{enumerate}
  \item اشتق الكثافة المتوسطة \(\rho_T\) من الحد الرئيس في صيغة ريمان--فون مانغولت.
  \item تحقق أن متوسط التباعد بعد الضرب في \(\log T/(2\pi)\) يقترب من واحد.
  \item أثبت من زوج فورييه \(\widehat K=\mathbf 1_{[-1/2,1/2]}\) أن \(\widehat{K^2}(\alpha)=(1-|\alpha|)_+\).
  \item بين لماذا يحتوي مجموع \(F(\alpha,T)\) على القطر، بينما تستبعده صيغة كثافة الأزواج غير القطرية.
  \item اشرح بمثال منطقي لماذا لا يحدد الارتباط الثنائي قانون أقرب جار.
  \item قارن بين تصنيف مبرهنة GUE المصفوفية وحدسية GUE لأصفار زيتا.
\end{enumerate}
